\section{Metodologia}

Este capítulo apresenta o desenho metodológico da pesquisa e distingue os procedimentos já executados daqueles que ainda integram o protocolo experimental. Até esta etapa, foram concluídas a constituição, o versionamento e a análise exploratória do \textit{corpus} de discursos. Permanecem planejadas a construção efetiva das unidades de recuperação, a elaboração do conjunto de perguntas e dos julgamentos de relevância, a comparação entre estratégias de busca e, em etapa posterior, a avaliação das respostas geradas. Essa distinção temporal evita tratar estimativas de preparação dos dados como resultados de recuperação ou geração.

\subsection{Natureza e abordagem da pesquisa}

A pesquisa possui natureza aplicada, exploratória e descritiva. É aplicada porque se volta a um problema concreto da administração pública: a necessidade de avaliar, com critérios explícitos, sistemas RAG voltados à consulta verificável de acervos legislativos extensos, tomando como caso empírico discursos parlamentares da 56\textordfeminine{} Legislatura do Senado Federal. É exploratória porque investiga a organização e a aplicação de um protocolo de avaliação de sistemas RAG, considerando a articulação entre recuperação documental, geração de respostas e verificabilidade das fontes. É descritiva porque já caracteriza o \textit{corpus} e, nas etapas experimentais, caracterizará os resultados da aplicação do protocolo segundo dimensões e critérios previamente definidos. Essa classificação segue a orientação metodológica de explicitar a natureza da pesquisa em função do problema investigado e do tipo de evidência necessário para respondê-lo \parencite{gil2022projetosPesquisa}.

Quanto ao método de abordagem, a pesquisa adota orientação predominantemente hipotético-dedutiva, pois parte de hipóteses de trabalho e de critérios avaliativos previamente definidos, mas incorporará componente indutivo ao identificar padrões de erro emergentes e formular recomendações a partir dos resultados observados. Essas hipóteses derivam principalmente da literatura sobre RAG, recuperação da informação e avaliação de sistemas generativos e são contextualizadas pela literatura sobre informação legislativa em contexto público; nas etapas experimentais, serão confrontadas com evidências empíricas do caso estudado. Ao mesmo tempo, o exame dos casos críticos poderá revelar padrões não previstos inicialmente, que serão usados para refinar recomendações metodológicas e institucionais. Essa combinação é adequada porque a pesquisa confronta respostas provisórias ao problema com evidências observáveis, mas também admite aprendizado a partir da observação sistemática dos resultados \parencite{lakatosMarconi2021metodologia}. Essa orientação não implica teste causal clássico entre variáveis isoladas; as hipóteses serão examinadas por meio de critérios definidos no protocolo avaliativo e aplicados ao sistema RAG, às respostas geradas e ao \textit{corpus} estudado.

A pesquisa adota abordagem multimétodos. A dimensão quantitativa já foi empregada na caracterização descritiva do \textit{corpus} e será ampliada pela mensuração da recuperação documental, pela avaliação estruturada da geração e, quando aplicável, por medidas de concordância entre julgamentos humanos e automatizados. A dimensão qualitativa já orientou a auditoria de casos documentais atípicos e estará presente na análise de erros, na interpretação das divergências entre avaliadores humanos e \textit{LLM-as-a-judge}, na interpretação institucional dos resultados e na leitura contextualizada de respostas representativas. O uso combinado de estratégias quantitativas e qualitativas é pertinente porque o fenômeno investigado envolve tanto desempenho mensurável quanto adequação institucional, qualidade argumentativa, interpretação documental e riscos de uso \parencite{creswellPlanoClark2018mixedMethods}.

Como delineamento, a dissertação combina pesquisa bibliográfica, pesquisa documental, estudo de caso aplicado e aplicação controlada do protocolo avaliativo. A pesquisa bibliográfica sistematiza literatura sobre RAG, avaliação de sistemas generativos, recuperação da informação, informação legislativa e uso institucional de IA. A pesquisa documental constituiu e auditou a base de discursos e metadados legislativos públicos. O estudo de caso concentra a análise em um sistema RAG aplicado aos discursos da 56\textordfeminine{} Legislatura do Senado Federal. A aplicação controlada do protocolo será usada para comparar recuperação documental, geração de respostas e verificabilidade das fontes sob critérios previamente definidos, sem pretensão de generalizar os resultados para todas as arquiteturas ou todos os domínios de RAG.

O enquadramento geral dialoga com a \textit{Design Science Research}, pois a dissertação organiza e aplicará um protocolo avaliativo como artefato metodológico voltado a enfrentar um problema relevante de informação pública, tendo um sistema RAG como caso empírico de aplicação \parencite{hevnerEtAl2004designScience,peffersEtAl2007dsrm}.

\subsection{Fundamentação empírica: dados e técnicas de coleta}

O \textit{corpus} documental foi constituído com dados públicos e secundários de pronunciamentos da 56\textordfeminine{} Legislatura do Senado Federal. A coleta consolidada abrange o período de 1 de fevereiro de 2019 a 10 de janeiro de 2023 e produziu 15.729 registros, cada qual identificado por um código de pronunciamento único. A versão utilizada na análise foi publicada como revisão \texttt{v1.1.1} e fixada pelo resumo SHA-256 \path{e09cfc4793e5394be440906320c7d3008cda5a52b90bbc85bf47446362406af1}. O uso do resumo criptográfico permite verificar a identidade binária do arquivo analisado, embora não substitua a auditoria de qualidade de seu conteúdo.

A constituição do \textit{corpus} preservou o esquema canônico empregado na prova de conceito e os campos adicionais retornados pela fonte oficial. Entre os metadados disponíveis estão data, autor, partido, unidade federativa, tipo de uso da palavra, resumo, indexação e URL do pronunciamento original. Os campos internos de início e fim da janela de coleta registram os intervalos mensais efetivamente consultados, sem alteração artificial para reproduzir a coleta anterior. Como se trata da organização de documentos públicos, não houve coleta direta de informações com participantes humanos nessa etapa.

As técnicas de coleta e produção de evidências distribuem-se em três frentes. A frente bibliográfica e técnico-normativa sistematiza artigos, \textit{surveys}, \textit{benchmarks}, relatórios institucionais e diretrizes sobre RAG, recuperação da informação, avaliação de sistemas generativos, informação legislativa e uso institucional de IA. A frente documental, já executada quanto à obtenção e ao congelamento da base, abrange coleta, consolidação, controle de versão e auditoria do \textit{corpus}. A frente avaliativa permanece em desenvolvimento e compreenderá o registro das unidades recuperadas, das respostas geradas e dos julgamentos humanos e automatizados definidos pelo protocolo.

\subsubsection{Análise exploratória e prontidão do corpus}

A análise exploratória foi executada no \textit{notebook} 01. O procedimento auditou esquema, tipos, unicidade dos códigos, validade das datas, valores ausentes, cobertura textual, distribuições temporal e institucional, extensão dos documentos, resíduos de HTML e URLs, caracteres problemáticos, duplicidade textual exata e aproximada, frequência lexical, mudança lexical anual, cobertura de metadados e concentração de representação. A análise de duplicidade aproximada empregou, em amostra reprodutível e controlada, assinaturas MinHash, candidatos por \textit{Locality-Sensitive Hashing}, similaridade de Jaccard sobre \textit{shingles} e similaridade TF-IDF. Esses procedimentos caracterizam riscos do \textit{corpus}; não avaliam a qualidade de um recuperador ou de um modelo gerador.

Dos 15.729 registros, 15.039 possuem texto integral, 687 dependem de resumo como fonte substituta e três não contêm texto nem resumo utilizável. Assim, 15.726 registros, ou 99,98\%, têm algum conteúdo indexável, mas texto integral e resumo não serão tratados como evidências equivalentes. A proveniência deverá acompanhar cada unidade recuperável para permitir que a avaliação distinga a fonte documental utilizada. Também foram identificadas 28 cópias textuais exatas excedentes, 1.291 documentos com menos de 50 palavras, sete documentos com mais de 10 mil palavras e cobertura de 86,81\% para partido e unidade federativa entre os documentos indexáveis. Esses achados recomendam inspeção antes de deduplicação, tratamento explícito de documentos curtos e longos e cautela no uso de filtros por metadados incompletos.

A estimativa de segmentação não criou unidades de recuperação reais. Com aproximação de 512 tokens por trecho e sobreposição de 64 tokens, o \textit{notebook} estimou 45.971 trechos, mediana de dois trechos para textos integrais e máximo aproximado de 60. Configurações de 256, 768 e 1.024 tokens também foram simuladas para dimensionar o índice. Esses valores servem ao planejamento; a segmentação efetiva deverá usar o tokenizador do modelo escolhido, respeitar limites semânticos e ser comparada empiricamente no protocolo de recuperação.

O quadro de prontidão adotou limiares heurísticos, e não critérios de certificação. O \textit{corpus} atendeu aos limiares exploratórios de indexabilidade, cobertura de texto integral, baixa proporção de registros sem conteúdo, baixa duplicidade exata excedente e proporção de documentos curtos. A cobertura mínima entre os metadados candidatos a filtro ficou abaixo do limiar heurístico de 90\%, em razão dos 86,81\% observados para partido e unidade federativa. Portanto, a conclusão desta etapa limita-se a considerar os dados aptos a avançar para experimentos controlados, com as ressalvas documentadas.

\subsubsection{Etapas ainda planejadas para o protocolo experimental}

O próximo ciclo experimental será iniciado no \textit{notebook} 02. Essa etapa construirá um conjunto de perguntas estratificado; associará a cada pergunta respondível os pronunciamentos ou trechos relevantes; produzirá efetivamente as unidades de recuperação; e comparará tamanhos e sobreposições de segmentação. Serão implementadas busca lexical com BM25, busca vetorial e busca híbrida, preservados os mesmos dados, perguntas e julgamentos de relevância para tornar a comparação controlada.

A bateria contemplará perguntas factuais, foco em autor, desambiguação, evidência negativa, comparação entre autores, comparação temporal, perguntas numéricas, síntese temática, múltiplas etapas, autoria cruzada, controle de escopo e estresse de recuperação. Essas categorias representam famílias recorrentes em avaliações de RAG e riscos específicos do domínio legislativo, como dispersão temporal, multiplicidade de autores, variação terminológica e perguntas não respondíveis \parencite{zhuEtAl2025rageval,saadFalconEtAl2024ares,esEtAl2023ragas,pipitoneAlami2024legalbenchRag,dahlEtAl2024largeLegalFictions}. Sua distribuição será estratificada por período, tamanho do documento e tipo de fonte e, quando a cobertura permitir, por partido e unidade federativa, sem interpretar ausência de metadados como categoria substantiva de representação.

O conjunto de referência será produzido por curadoria documental, com preservação do vínculo entre trecho, pronunciamento e URL oficial. Perguntas não respondíveis terão essa condição registrada explicitamente. Também serão documentadas ambiguidades, decisões de relevância e eventuais revisões do conjunto. Somente depois desse conjunto será possível calcular métricas de recuperação sem confundir similaridade automática com relevância de referência.

Após a estabilização da recuperação, o protocolo avançará para a geração de respostas e para a comparação entre rubrica humana e julgamento automatizado. Essa fase permanecerá analiticamente separada da recuperação: uma resposta inadequada poderá decorrer da ausência de evidência pertinente, do uso incorreto de evidência disponível ou da extrapolação promovida pelo modelo gerador. Serão registrados estratégia de segmentação, modelos, \textit{embeddings}, índices, parâmetros de busca, valor de \textit{top-k}, reranqueamento, \textit{prompt}, modelo gerador, temperatura, \textit{top-p}, versões e datas de execução. Nesse vocabulário, \textit{top-k} designa a quantidade de resultados documentais considerada na recuperação, enquanto \textit{top-p} controla a amostragem probabilística de tokens na geração.

\subsection{Técnicas de análise dos dados e matriz de opções metodológicas}

A análise experimental será organizada em quatro camadas, posteriores à caracterização do \textit{corpus}. A primeira avaliará a recuperação documental por meio de métricas de recuperação ranqueada consolidadas na literatura de recuperação da informação, como \textit{Recall@k}, \textit{Precision@k}, MRR e nDCG \parencite{manningEtAl2008informationRetrieval}, também empregadas em \textit{benchmarks} contemporâneos como o BEIR \parencite{thakurEtAl2021beir}. O \textit{notebook} 02 calculará \textit{Recall@1}, \textit{Recall@5} e \textit{Recall@10}, além de MRR e nDCG@k; \textit{Precision@k} e \textit{Hit@k} serão mantidas quando compatíveis com a granularidade dos julgamentos de relevância. A comparação abrangerá BM25, recuperação vetorial e recuperação híbrida sob o mesmo conjunto de referência.


Para fins de avaliação da recuperação, a unidade principal de análise será o trecho recuperável indexado no sistema, preservada sua vinculação ao discurso original, aos metadados e à URL da fonte. Quando a evidência esperada corresponder a um discurso inteiro, serão considerados relevantes os trechos desse discurso que contenham informação suficiente para responder à pergunta.

\textit{Recall@k} indicará a proporção das evidências esperadas presente entre os primeiros resultados; \textit{Precision@k} indicará a proporção de itens relevantes entre os recuperados; MRR observará a posição da primeira evidência relevante; nDCG considerará a ordenação das evidências segundo seus graus de relevância; e \textit{Hit@k} indicará se ao menos uma evidência suficiente foi recuperada. Além dos resultados agregados, o desempenho será estratificado por ano, tamanho do documento, tipo de fonte e demais metadados cuja cobertura permita comparação válida. A diversidade dos resultados e os erros de recuperação serão examinados separadamente para identificar redundância, concentração e falhas sistemáticas.

O uso de \textit{Precision@k} e nDCG dependerá da granularidade do conjunto de referência e da possibilidade de classificar os itens recuperados como relevantes, parcialmente relevantes ou irrelevantes. Caso o conjunto de referência registre apenas evidências suficientes por pergunta, a análise priorizará \textit{Hit@k}, \textit{Recall@k} e MRR, mantendo \textit{Precision@k} e nDCG apenas quando houver julgamento de relevância suficientemente detalhado para sua aplicação.

A segunda camada avaliará a geração de respostas por rubrica, considerando correção factual, aderência ao contexto recuperado, fidelidade às fontes, precisão de autoria, explicitação de insuficiência documental, completude proporcional e ausência de alucinação. A terceira camada comparará avaliação humana e avaliação baseada no paradigma \textit{LLM-as-a-judge}, observando concordâncias, divergências, falsos positivos de qualidade e falsos negativos. A quarta camada analisará qualitativamente erros e casos críticos, buscando identificar padrões como falhas de recuperação, uso inadequado de fontes, atribuição incorreta, extrapolação do contexto recuperado, resposta sem evidência suficiente e ausência de explicitação de insuficiência documental.

Essas camadas operacionalizam as hipóteses do estudo e mantêm correspondência direta com os objetivos específicos. A sistematização da literatura orientará a seleção das métricas, rubricas e categorias de pergunta. A caracterização do acervo delimitará o \textit{corpus}, os metadados e os riscos documentais relevantes. A bateria de perguntas e o conjunto de referência permitirão examinar recuperação documental, geração de respostas e verificabilidade das fontes em condições controladas. A comparação entre avaliação humana e avaliação baseada no paradigma \textit{LLM-as-a-judge} permitirá examinar o alcance do julgamento automatizado. Por fim, a análise qualitativa dos erros e casos críticos sustentará a discussão sobre limites, riscos e condições de interpretação dos resultados em contexto institucional. Assim, a verificabilidade será examinada pela relação entre evidências recuperadas, fontes citadas e afirmações geradas; a variação por tipo de pergunta será analisada por desempenho agregado e padrões de erro; o alcance da avaliação baseada no paradigma \textit{LLM-as-a-judge} será avaliado pela comparação com julgamentos humanos independentes; e as condições de interpretação e uso cauteloso serão discutidas a partir dos resultados técnicos, dos casos críticos e dos requisitos institucionais identificados na literatura.

O protocolo avaliativo incluirá: aplicação a um sistema RAG construído sobre o \textit{corpus} documental; bateria de perguntas distribuída por categorias; conjunto de referência com evidências documentais associadas às perguntas respondíveis; registro das evidências recuperadas e das respostas geradas; avaliação humana independente de uma amostra de respostas; e avaliação auxiliar baseada no paradigma \textit{LLM-as-a-judge}. Esses componentes serão utilizados para examinar recuperação documental, geração de respostas e verificabilidade das fontes, sem pretensão de comparar arquiteturas RAG de forma abrangente.

A avaliação humana será organizada por rubrica explícita, preferencialmente aplicada por pelo menos dois avaliadores independentes com familiaridade mínima com documentação legislativa ou com avaliação de sistemas de informação. Antes da anotação principal, será realizado piloto com amostra reduzida de perguntas para calibrar instruções, exemplos positivos, exemplos negativos e casos limítrofes. Sempre que viável, os avaliadores não terão acesso à identidade do sistema que gerou cada resposta. Divergências serão registradas e, quando necessário, dirimidas por revisão conjunta ou por avaliador adicional. A concordância entre julgamentos será descrita por percentual de acordo e pela análise das divergências observadas. A interpretação dos resultados será cautelosa e proporcional ao desenho amostral, evitando apresentar médias agregadas ou percentuais de concordância como prova isolada de qualidade.

Os avaliadores serão informados sobre a finalidade da avaliação, o uso dos resultados e a forma de registro das respostas, preservando-se sua identificação individual na apresentação dos resultados.

Para o julgamento automatizado, serão registrados modelo, versão, \textit{prompt}, rubrica, temperatura, formato de entrada e data de execução. Sempre que viável, o julgamento será repetido ou validado em amostra piloto para observar estabilidade. Os resultados do \textit{LLM-as-a-judge} serão tratados como auxiliares e comparados à avaliação humana, sem função de certificação autônoma.

\subsection{Recursos, viabilidade e transparência no uso de IA}

A viabilidade da pesquisa depende de quatro condições principais. A primeira é a disponibilidade do \textit{corpus}, composto por dados públicos e secundários, o que reduz riscos de acesso e dispensa coleta direta de dados pessoais sensíveis junto a participantes humanos para constituição da base documental. A segunda é a infraestrutura computacional para preparação textual, indexação vetorial, execução do sistema RAG, armazenamento de resultados e registro de versões; essa infraestrutura poderá ser composta por ferramentas de código aberto e, quando necessário, por serviços de API para geração ou julgamento automatizado. A terceira é a disponibilidade de avaliadores humanos para aplicar rubrica explícita a uma amostra das respostas; por isso, a avaliação humana será tratada como etapa amostral, com piloto prévio e escopo ajustável ao número efetivo de avaliadores. A quarta é o controle de custos e tempo de execução, que será enfrentado pela delimitação do protocolo avaliativo, pela aplicação a um \textit{corpus} documental específico e pela realização da avaliação humana em amostra de respostas.

O uso de IA generativa na pesquisa será tratado de forma transparente. Modelos de linguagem poderão ser utilizados como componentes do sistema RAG avaliado, como instrumentos auxiliares de julgamento automatizado e, como apoio editorial à revisão de texto, sempre sob responsabilidade do pesquisador. Quando integrarem a coleta ou a análise dos dados, serão registrados modelo, versão ou provedor disponível, configuração relevante, rubrica, prompt ou instruções de julgamento e data de execução. As saídas automatizadas não serão tratadas como evidência autossuficiente: deverão ser confrontadas com fontes documentais, métricas de recuperação, rubricas e avaliação humana nos limites definidos pelo protocolo.

A seguir apresenta-se a matriz preliminar de opções metodológicas que a pesquisa empregará:

\begin{landscape}
\begingroup
\scriptsize
\setlength{\tabcolsep}{1.5pt}
\begin{longtable}{@{}>{\raggedright\arraybackslash}p{0.13\linewidth}>{\raggedright\arraybackslash}p{0.095\linewidth}>{\raggedright\arraybackslash}p{0.075\linewidth}>{\raggedright\arraybackslash}p{0.115\linewidth}>{\raggedright\arraybackslash}p{0.095\linewidth}>{\raggedright\arraybackslash}p{0.09\linewidth}>{\raggedright\arraybackslash}p{0.105\linewidth}>{\raggedright\arraybackslash}p{0.125\linewidth}>{\raggedright\arraybackslash}p{0.085\linewidth}@{}}
\caption{Matriz preliminar de opções metodológicas}
\label{tab:matriz-metodologica}\\
\toprule
\textbf{Objetivo específico} & \textbf{Abordagem} & \textbf{Natureza} & \textbf{Dados coletados} & \textbf{Público-alvo} & \textbf{Amostragem} & \textbf{Técnicas de coleta} & \textbf{Técnicas de análise} & \textbf{Apresentação} \\
\midrule
\endfirsthead
\toprule
\textbf{Objetivo específico} & \textbf{Abordagem} & \textbf{Natureza} & \textbf{Dados coletados} & \textbf{Público-alvo} & \textbf{Amostragem} & \textbf{Técnicas de coleta} & \textbf{Técnicas de análise} & \textbf{Apresentação} \\
\midrule
\endhead
Sistematizar literatura sobre avaliação de RAG & Qualitativa & Exploratória & Artigos, relatórios, diretrizes e \textit{benchmarks} & Pesquisadores e gestores públicos & Intencional por pertinência temática & Pesquisa bibliográfica e documental & Síntese crítica por temas e lacunas & Texto e quadro-síntese \\
Selecionar técnicas de avaliação aplicáveis ao caso & Multimétodos & Exploratória & Métricas, rubricas, critérios e protocolos & Gestores, equipes técnicas e avaliadores & Seleção por adequação ao problema & Fichamento e análise de protocolos & Comparação metodológica e análise de viabilidade & Protocolo de avaliação \\
Caracterizar o acervo legislativo & Qualitativa & Descritiva & Discursos, metadados e documentação do \textit{corpus} & Usuários de informação legislativa & \textit{Corpus} da 56\textordfeminine{} Legislatura & Pesquisa documental e registro de versões & Caracterização documental e análise de metadados & Descrição do \textit{corpus} \\
Construir bateria e conjunto de referência & Multimétodos & Exploratória & Perguntas, respostas esperadas e evidências & Avaliadores e pesquisadores & Perguntas balanceadas por categoria & Curadoria documental e anotação das evidências & Categorização, revisão e validação & Bateria e quadro de evidências \\
Avaliar recuperação documental & Quantitativa & Descritiva & Trechos recuperados e evidências esperadas & Equipes técnicas e pesquisadores & Perguntas da bateria & Execuções controladas do protocolo avaliativo & Hit@k, Recall@k, Precision@k, MRR e nDCG & Tabelas e gráficos \\
Avaliar geração e julgamentos & Multimétodos & Descritiva & Respostas, notas humanas e notas de avaliação baseada no paradigma \textit{LLM-as-a-judge} & Avaliadores humanos e gestores & Amostra de respostas, preservando representação das categorias de pergunta & Rubrica humana, piloto e avaliação baseada no paradigma \textit{LLM-as-a-judge} & Estatísticas, concordância e análise de erros & Tabelas, exemplos e discussão \\
Analisar limites e condições de interpretação dos resultados & Qualitativa & Descritiva & Resultados, erros, divergências avaliativas e casos críticos & Gestores públicos, equipes técnicas e pesquisadores & Casos críticos e achados do protocolo & Síntese dos resultados, análise de erros e seleção de casos representativos & Análise temática de erros, limites e implicações institucionais & Discussão interpretativa e recomendações de uso cauteloso \\
\bottomrule
\end{longtable}

\noindent Fonte: elaborado pelo autor.
\endgroup
\end{landscape}
